%%% ВЫБОР УСТРОЙСТВА %%%

\IfStrEqCase{\devicecode}{%
{ЮТКБ.656122.611-06.01 РЭ}{%ДЗТ
\def\fileExt{t}
\def\AbzazOne{
\item
Для исключения пуска пожаротушения не отключенного от сети трансформатора обеспечиваются следующие виды контроля:
\begin{itemize}
\item отсутствие тока на стороне ВН;
\item отсутствие тока на сторонах СН и НН;
\item отсутствие напряжения на сторонах СН и НН.
\end{itemize}
}
}%
{ЮТКБ.656122.611-06.02 РЭ}{%ДЗАТ
\def\fileExt{at}
\def\AbzazOne{
\item
Для исключения пуска пожаротушения не отключенного от сети трансформатора обеспечиваются следующие виды контроля:
\begin{itemize}
\item отсутствие тока на стороне ВН;
\item отсутствие тока на стороне СН;
\item отсутствие напряжения на стороне НН.
\end{itemize}
}
}%
}
[
\def\fileExt{??}
\def\AbzazOne{??}
]

\needspace{2\baselineskip}
\color{unidarkgreen}{\subsection{Контроль отсутствия напряжения}}
\color{black}

\begin{enumerate}

\AbzazOne

\item
Алгоритм работы КОН приведен на рисунке \ref{pic:KON}. Параметры для настройки КОН приведены в таблице \ref{app_set:kon}.

\medskip
\begin{figure}[H]
\centering
\includegraphics[width=1\textwidth,height=1\textheight,keepaspectratio]{\fbpath/06.02 ДЗАТ/041.1201 КОН/img/KON_\fileExt.pdf}
\captionof{figure}{Функциональная схема логики КОН}\label{pic:KON}
\end{figure}


\end{enumerate}